\subsection{Polynomial obstructions and adversary synthesis certificates}\label{ian:polynomial-adversary-certificates}

Status: proved: the global degree and amplitude conditions, the coefficient obstruction \eqref{ian:pac-pm-display} and its arbitrary-depth form \eqref{ian:pac-pm-tail}, product-domination exclusion, conditional adversary-to-sector constructions, an exact query-budget synthesis certificate, and the character-code and Simon checks. Refuted: the claims that completion weights alone imply a quantum query bound depending only on their cost (Proposition~\ref{ian:pac-completion-not-synthesis}), and that a large adversary lower bound alone certifies a proposed realization (Remark~\ref{ian:pac-realization-gap}). General adversary synthesis supplies approximate decision algorithms; matching another prescribed acceptance function remains a separate requirement. Literature imports, each marked where used: the polynomial method of Beals--Buhrman--Cleve--Mosca--de Wolf (Lemmas~4.1--4.2), the query-state semidefinite characterization of Barnum--Saks--Szegedy, the general adversary lower bound of H{\o}yer--Lee--{\v S}palek (Theorem~2), Reichardt's dual adversary and synthesis, and Lee--Mittal--Reichardt--{\v S}palek--Szegedy (Theorem~1.1) for general alphabets; the Bernstein--Vazirani and Simon procedures are cited for orientation, and the calculations used are proved here. Theorem~\ref{ian:pac-product-exclusion} and the first assertion of Proposition~\ref{ian:pac-completion-weights} rest on \path{foundations/product_alphabets.tex}, which imports inequality (BC) of Huffstutler--Miloschewsky--Podder. Repository inputs are Lemma~\ref{lem:resolution-symbol-adversary} of \path{foundations/product_alphabets.tex}, Lemma~\ref{ian:degree-threshold-certificates} of \path{ian/22_influence_form_uniform_aa.tex}, and Proposition~\ref{ian:hard-sector-forward} and Corollary~\ref{ian:hard-sector-explicit} of \path{ian/13_hard_sector_definitions_forward.tex}.

Throughout, \(\Omega_n=\{-1,1\}^n\), and a coordinate query is
\[
 O_x|i,b\rangle
   =|i,b\oplus(1-x_i)/2\rangle.
\]
Input-independent unitaries, workspace and classical postprocessing are free. No sampler, promise-membership oracle or hidden-label oracle is included. Measurements and randomness may be purified and deferred.

We use Definition~\ref{ian:hard-sector-definition} in
\path{ian/13_hard_sector_definitions_forward.tex}.
Thus, for \(Q:\Omega_n\to[0,1]\),
\[
 \operatorname{Hard}(Q,\nu;t,a,\gamma)
 \quad\Longleftrightarrow\quad
 \operatorname{Var}_\nu Q\ge a,\qquad
 \sup_{F\in\mathcal F_t}
       |\operatorname{Cov}_\nu(Q,F)|\le\gamma,
 \qquad 0<\gamma<a\le\tfrac14.
\]
Here \(\mathcal F_t\) contains all deterministic adaptive coordinate
trees of worst-case depth at most \(t\), with outputs in \([-1,1]\).
A smooth sector of ambient weight \(w\) satisfies \(w\nu\le\mu\).
For a set sector, \(\nu=\mu(\cdot\mid H)\) and \(\mu(H)\ge w\).
When discussing a restriction to \(H\), we assume \(\nu(H)=1\).
As in \path{ian/13_hard_sector_definitions_forward.tex}, \(\mathcal Q_{n,T}\) is used with \(T\ge1\).

\begin{lemma}[Global polynomial and amplitude conditions]
\label{ian:pac-polynomial-sos}
If \(Q\in\mathcal Q_{n,T}\), its unique multilinear representation
on the whole cube has degree at most \(\min\{2T,n\}\), and
\(0\le Q(x)\le1\) for every \(x\in\Omega_n\).
There are finite families of real multilinear polynomials
\(u_j,v_k\), each of degree at most \(T\), such that
\[
 Q(x)=\sum_j u_j(x)^2,\qquad
 1-Q(x)=\sum_k v_k(x)^2
 \quad(x\in\Omega_n).
 \tag{32.1}\label{ian:pac-sos}
\]
These are identities on the cube, equivalently modulo
\(x_i^2-1\).

More explicitly, put
\[
 \mathcal S_T=\{S\subseteq[n]:|S|\le T\},
 \qquad z_T(x)=(\chi_S(x))_{S\in\mathcal S_T},
 \qquad \chi_S(x)=\prod_{i\in S}x_i.
\]
There are real positive semidefinite matrices \(M_1,M_0\), indexed
by \(\mathcal S_T\), with
\[
 Q=z_T^{\mathsf T}M_1z_T,\qquad
 1-Q=z_T^{\mathsf T}M_0z_T
 \quad\hbox{on }\Omega_n.
\]
Consequently, for every \(U\subseteq[n]\),
\[
 \widehat Q(U)
   =\sum_{\substack{S,R\in\mathcal S_T\\S\triangle R=U}}
       (M_1)_{S,R},
 \qquad
 \sum_{\substack{S,R\in\mathcal S_T\\S\triangle R=U}}
       (M_0+M_1)_{S,R}
   =\mathbf1_{\{U=\varnothing\}}.
 \tag{32.2}\label{ian:pac-sos-coefficients}
\]
In particular, \(\operatorname{tr}(M_0+M_1)=1\).
\end{lemma}

\begin{proof}
In the response basis
\(|+\rangle,|-\rangle\), the oracle acts as the identity on
\(|i,+\rangle\) and multiplication by \(x_i\) on
\(|i,-\rangle\). Initially all amplitudes are constant.
Input-independent unitaries take linear combinations of amplitudes,
and each query increases their degree by at most one.
Reducing \(x_i^2=1\) preserves their values on the cube and cannot
increase degree. Every final amplitude therefore has a multilinear
representation of degree at most \(T\).

After adjoining an output register, acceptance and rejection are
orthogonal sets of basis states. Their probabilities are the
respective sums of squared absolute amplitudes. Splitting each
amplitude into its real and imaginary parts proves
\eqref{ian:pac-sos}. Normalization gives their sum equal to one.
Multilinearizing the squared expressions proves the degree bound.
Writing each real amplitude in the character basis gives the
positive semidefinite matrices. Finally,
\(\chi_S\chi_R=\chi_{S\triangle R}\), which proves
\eqref{ian:pac-sos-coefficients}.

The amplitude and probability statements are the polynomial method
of Beals--Buhrman--Cleve--Mosca--de Wolf (from literature:
\emph{Quantum Lower Bounds by Polynomials}, J.~ACM 48(4),
778--797, 2001; Lemmas 4.1--4.2 in arXiv:quant-ph/9802049v3).
\end{proof}

\begin{proposition}[The coefficient obstruction]
\label{ian:pac-pm}
Suppose \(Q\in\mathcal Q_{n,T}\) and
\(\operatorname{Hard}(Q,\nu;t,a,\gamma)\).
Write \(d=\min\{2T,n\}\) and
\[
 Q(x)=\sum_{|S|\le d}\widehat Q(S)\chi_S(x),
 \qquad
 \widehat Q(S)=\mathbb E_{U_n}[Q\chi_S].
\]
If \(t\ge d\), then
\[
 a\le
 \gamma\sum_{1\le |S|\le d}|\widehat Q(S)|.
 \tag{32.3}\label{ian:pac-pm-display}
\]
Since \(d\le2T\), this applies in particular whenever \(t\ge2T\); the choice \(d=\min\{2T,n\}\) refines that form of the depth hypothesis.

For arbitrary \(t\), the valid replacement is
\[
 a\le
 \gamma\sum_{1\le |S|\le\min\{t,d\}}|\widehat Q(S)|
 +\frac12\sum_{t<|S|\le d}|\widehat Q(S)|.
 \tag{32.4}\label{ian:pac-pm-tail}
\]
If \(t\ge d\), global boundedness also gives
\[
 \frac a\gamma
 \le\sum_{1\le|S|\le d}|\widehat Q(S)|
 \le\frac12
       \sqrt{\sum_{j=1}^d\binom nj}.
\]
\end{proposition}

\begin{proof}
Each \(\chi_S\) is a deterministic depth-\(|S|\) coordinate tree
with outputs in \(\{-1,1\}\). Thus
\[
 |\operatorname{Cov}_\nu(Q,\chi_S)|\le\gamma
 \qquad(|S|\le t).
\]
Taking covariance with the polynomial identity gives
\[
 \operatorname{Var}_\nu Q
   =\sum_{\varnothing\ne S,\ |S|\le d}
        \widehat Q(S)\operatorname{Cov}_\nu(Q,\chi_S).
\]
This proves \eqref{ian:pac-pm-display} when \(t\ge d\).
For the remaining terms, Cauchy--Schwarz and boundedness give
\[
 |\operatorname{Cov}_\nu(Q,\chi_S)|
 \le\sqrt{\operatorname{Var}_\nu Q\,
                 \operatorname{Var}_\nu\chi_S}\le\tfrac12,
\]
proving \eqref{ian:pac-pm-tail}.
Parseval under the uniform cube law gives
\[
 \sum_{S\ne\varnothing}\widehat Q(S)^2
   =\operatorname{Var}_{U_n}Q\le\tfrac14,
\]
and the last assertion follows by Cauchy--Schwarz.
\end{proof}

\begin{remark}[Coefficients on a sector]
\label{ian:pac-sector-coefficients}
The coefficients in \eqref{ian:pac-pm-display} belong to a polynomial on the physical
cube. They are not Fourier coefficients defined by the possibly
dependent sector law. More generally, the same proof works with
any polynomial \(p\) of degree at most \(d\) agreeing with \(Q\)
on \(\operatorname{supp}\nu\), using
\(\operatorname{Var}_\nu Q=\operatorname{Cov}_\nu(Q,p)\).
A promise restriction need not have a unique polynomial extension.

A zero-query algorithm has constant acceptance, so it has no hard sector.
Nor can hardness hold at \(t\ge n\): the tree
\(F=2Q-1\) then gives covariance
\(2\operatorname{Var}_\nu Q\ge2a>\gamma\).
\end{remark}

\begin{theorem}[Product-domination exclusion]
\label{ian:pac-product-exclusion}
There is an absolute constant \(C_{\mathrm P}\ge1\) with the
following property. Let \(T\ge1\), let
\(Q\in\mathcal Q_{n,T}\), and suppose
\[
 \operatorname{Hard}(Q,\nu;t,a,\gamma),
 \qquad \nu\le K\pi,
\]
where \(\pi\) is a Boolean product law and \(K\ge1\).
Then
\[
 t<
 \left\lceil
 C_{\mathrm P}T^6
       \left(\frac{2K}{a-\gamma}\right)^{15}
 \right\rceil.
 \tag{32.5}
\]
Consequently, if a sector of ambient weight \(w\) also satisfies
\(\mu\le L\pi\), then
\[
 t<
 \left\lceil
 C_{\mathrm P}T^6
       \left(\frac{2L}{w(a-\gamma)}\right)^{15}
 \right\rceil.
\]
In particular, a sector satisfying \(w\nu\le\pi\) is excluded
at or beyond this bound with \(L=1\).
\end{theorem}

\begin{proof}
Corollary~\ref{cor:resolution-product-alphabet-aa} in
\path{foundations/product_alphabets.tex}, restricted to Boolean
alphabets, supplies a deterministic simulator \(D\) with
\[
 \mathbb E_\pi|D-Q|\le e,\qquad
 \operatorname{depth}(D)
   \le\left\lceil C_{\mathrm P}T^6e^{-15}\right\rceil.
\]
Choose \(e=(a-\gamma)/(2K)\), which lies in \((0,1)\).
Then \(\mathbb E_\nu|D-Q|\le Ke=(a-\gamma)/2\).
But Proposition~\ref{ian:hard-sector-forward}, applied with ambient
law \(\nu\) and weight one, gives
\[
 \mathbb E_\nu|D-Q|\ge a-\gamma
 \qquad(D\in\mathcal D_t).
\]
Thus the displayed simulator cannot have depth at most \(t\).
For the ambient formulation, \(w\nu\le\mu\le L\pi\), so take
\(K=L/w\).
\end{proof}

\begin{remark}[The accompanying influence condition]
\label{ian:pac-product-influence}
The same global polynomial is subject to the proved product
influence theorem in
\path{ian/22_influence_form_uniform_aa.tex}.
For every product law \(\pi\), define
\[
 \mathcal E_i^\pi(Q)
   =\mathbb E_\pi\operatorname{Var}(Q\mid X_{-i}).
\]
Corollary~\ref{ian:product-influence-aa-proved} gives
\[
 \max_i\mathcal E_i^\pi(Q)
 \ge c_\star
       \frac{\operatorname{Var}_\pi(Q)^{12}}{(2T)^{12}},
 \qquad c_\star=\frac9{2^{61}\pi^{12}}.
\]
The assertion holds after every positive-probability pinning.
For the uniform law it is also
Theorem~\ref{ian:uniform-influence-aa-proved}.
These assertions concern product laws and do not replace the
covariance requirement under a dependent sector law.
\end{remark}

\begin{proposition}[What the repository completion weights give]
\label{ian:pac-completion-weights}
Let \(f:S\to\{0,1\}\), with \(S\subseteq\Omega_n\).
Suppose nonnegative weights satisfy
\[
 \sum_i\lambda_i^x\le A,\qquad
 \sum_{i:x_i\ne y_i}\min(\lambda_i^x,\lambda_i^y)\ge1
       \quad(f(x)\ne f(y)).
 \tag{32.6}\label{ian:pac-cw}
\]
For every product law \(\pi\) and every \(\epsilon>0\), there is a
deterministic tree \(D:\Omega_n\to\{0,1\}\) with
\[
 \operatorname{depth}(D)
 \le\left\lceil128A^3\epsilon^{-2}\right\rceil,
 \qquad
 \pi\{x\in S:D(x)\ne f(x)\}\le\epsilon.
\]
This is unconditional error mass on the promise, not error under
\(\pi(\cdot\mid S)\).

If \(\nu(S)=1\), \(\nu\le K\pi\), and
\(\mathbb E_\nu|Q-f|\le\delta\), then this tree satisfies
\[
 \mathbb E_\nu|D-Q|\le K\epsilon+\delta.
\]
It therefore excludes
\(\operatorname{Hard}(Q,\nu;t,a,\gamma)\) whenever
\[
 t\ge\left\lceil128A^3\epsilon^{-2}\right\rceil,
 \qquad K\epsilon+\delta<a-\gamma.
 \tag{32.7}\label{ian:pac-completion-exclusion}
\]

Conversely, a \(q\)-query algorithm deciding \(f\) with error at
most \(1/3\) supplies weights \eqref{ian:pac-cw} of cost
\[
 A\le\frac{2q^2}{c_0^2},
 \qquad c_0=\frac{1-2\sqrt2/3}{2}>0.
\]
Also, if \(Q\in\mathcal Q_{n,T}\) and
\[
 f(x)=\mathbf1_{\{Q(x)\ge\theta\}},\qquad
 S=\{x:|Q(x)-\theta|\ge h\},\qquad h>0,
\]
then weights \eqref{ian:pac-cw} exist with
\[
 A\le\frac{2\pi^2T^2}{h}.
\]
\end{proposition}

\begin{proof}
The first assertion is
Theorem~\ref{thm:resolution-alphabet-completion} in
\path{foundations/product_alphabets.tex}.
Domination and the triangle inequality give the second assertion;
the sector forward bound gives
\eqref{ian:pac-completion-exclusion}.

For the algorithm-to-weights assertion, let \(p_{r,i}(x)\) be the
query probability at index \(i\) immediately before query \(r\), and put
\(a_i(x)=\sum_r p_{r,i}(x)\), so \(\sum_i a_i(x)=q\).
Lemma~\ref{lem:resolution-symbol-adversary} proves
\[
 \sum_{i:x_i\ne y_i}\sqrt{a_i(x)a_i(y)}\ge c_0
 \quad(f(x)\ne f(y)).
\]
Cauchy--Schwarz gives
\[
 c_0^2
 \le2q\sum_{i:x_i\ne y_i}\min(a_i(x),a_i(y)).
\]
Thus \(\lambda_i^x=(2q/c_0^2)a_i(x)\) works.

Finally, apply Lemma~\ref{ian:degree-threshold-certificates} in
\path{ian/22_influence_form_uniform_aa.tex}, with \(F=Q\),
degree bound \(2T\), threshold \(\theta\) and margin \(h\).
Its cost is \(\pi^2(2T)^2/(2h)\).
\end{proof}

\begin{proposition}[Completion weights alone do not synthesize algorithms]
\label{ian:pac-completion-not-synthesis}
There is no quantum query upper bound depending only on the cost
\(A\) in \eqref{ian:pac-cw}, even for Boolean input and output and
any fixed decision error below \(1/2\).
\end{proposition}

\begin{proof}
Fix a proposed query budget \(T\).
For arbitrarily large \(n\) divisible by four, there is a binary
code \(S\subseteq\Omega_n\) with minimum Hamming distance at least
\(n/4\) and
\[
 |S|>\sum_{j=0}^{2T}\binom nj.
\]
Indeed, greedy packing gives at least
\[
 \frac{2^n}{\sum_{j<n/4}\binom nj}
 \ge 2^{(1-H_2(1/4))n}
\]
codewords, whereas the right-hand side of the preceding inequality
is polynomial in \(n\) for fixed \(T\).
Here \(H_2(p)=-p\log_2p-(1-p)\log_2(1-p)\).
The binomial estimate follows by expanding
\((p+(1-p))^n\) at \(p=1/4\).

The restrictions to \(S\) of degree-at-most-\(2T\) polynomials
therefore have a nonzero real annihilator
\((c_x)_{x\in S}\):
\[
 \sum_{x\in S}c_xp(x)=0
 \quad\hbox{for every such polynomial }p.
\]
Since constants are included, \(\sum_xc_x=0\).
Set \(f(x)=\mathbf1_{\{c_x>0\}}\). Then
\[
 \sum_xc_xf(x)=\tfrac12\sum_x|c_x|.
\]
If a degree-at-most-\(2T\) polynomial \(p\) satisfied
\(|p(x)-f(x)|\le\delta<1/2\) on \(S\), we would have
\[
 \tfrac12\sum_x|c_x|
   =\left|\sum_xc_x(f(x)-p(x))\right|
   \le\delta\sum_x|c_x|,
\]
a contradiction. Lemma~\ref{ian:pac-polynomial-sos} consequently
excludes a \(T\)-query decision algorithm with that error.

Nevertheless \(\lambda_i^x=4/n\) satisfies
\eqref{ian:pac-cw}, for every labeling of \(S\), with cost \(A=4\).
\end{proof}

\begin{definition}[An exact certificate for a prescribed query budget]
\label{ian:pac-query-certificate}
Fix a finite promise \(S\subseteq\Omega_n\), a target
\(q:S\to[0,1]\), nonnegative tolerances \(\epsilon_x\), and an
integer \(T\ge1\).
Put \(x_0=1\), let \(J\) be the all-ones matrix on \(S\), and set
\[
 (C_i)_{x,y}=x_i y_i\qquad(0\le i\le n).
\]
A query certificate consists of Hermitian positive semidefinite matrices
\(B_i^r,E_0,E_1\), and Hermitian matrices \(G^r\), satisfying
\[
 G^0=J,
 \qquad
 \sum_{i=0}^n B_i^r=G^{r-1},
 \qquad
 G^r=\sum_{i=0}^n C_i\circ B_i^r
 \quad(1\le r\le T),
 \tag{32.8}\label{ian:pac-query-sdp}
\]
\[
 G^T=E_0+E_1,
 \qquad
 |(E_1)_{x,x}-q(x)|\le\epsilon_x
 \quad(x\in S).
 \tag{32.9}
\]
The product \(\circ\) is entrywise multiplication.
Positive semidefinite factorizations and the displayed equalities
and inequalities constitute a finite certificate.
Since \(J\) and every \(C_i\) are real, replacing every matrix of a certificate by its entrywise real part gives another certificate: the real part of a Hermitian positive semidefinite matrix is the average of the matrix and its entrywise conjugate, both positive semidefinite. Real symmetric certificates therefore suffice.
\end{definition}

\begin{theorem}[Exact synthesis from query-state feasibility]
\label{ian:pac-query-synthesis}
The certificate in Definition~\ref{ian:pac-query-certificate}
exists if and only if there is a globally defined acceptance
function \(Q\in\mathcal Q_{n,T}\) satisfying
\[
 |Q(x)-q(x)|\le\epsilon_x\qquad(x\in S).
\]
In particular, setting all tolerances to zero certifies exact
agreement on the promise. Taking \(S=\Omega_n\) certifies a
prescribed acceptance function on the whole cube.
\end{theorem}

\begin{proof}
This characterization is due to Barnum--Saks--Szegedy (from literature:
H.~Barnum, M.~Saks and M.~Szegedy, \emph{Quantum query complexity and semi-definite programming},
IEEE Conference on Computational Complexity (CCC 2003)); we give the proof in the present notation.

For necessity, purify an algorithm and use the response eigenbasis.
An algorithm making fewer than \(T\) queries is first padded with inert queries, using response \(|+\rangle\) at a fixed index; they act as the identity and leave the acceptance function unchanged.
The query space decomposes into an inert channel and channels
\(i=1,\ldots,n\) multiplied by \(x_i\).
Let \(B_i^r\) be the Gram matrix of the channel-\(i\) components
immediately before query \(r\). Their sum is the Gram matrix
before that query. Input-independent unitaries preserve Gram
matrices, while querying multiplies the channel-\(i\) Gram matrix
entrywise by \(C_i\). This gives
\eqref{ian:pac-query-sdp}. Acceptance and rejection components
give \(E_1,E_0\).

For sufficiency, factor
\[
 (B_i^r)_{x,y}=\langle b_{x,i}^r,b_{y,i}^r\rangle.
\]
The vectors
\[
 \sum_{i=0}^n |i\rangle|b_{x,i}^r\rangle
\]
have Gram matrix \(G^{r-1}\), and after the phase query they have
Gram matrix \(G^r\).
Two finite vector families with the same Gram matrix admit an
isometry mapping one family to the other: the map on their spans
is well-defined because every linear relation has squared norm
determined by that Gram matrix. After enlarging workspace, extend
each isometry to a unitary.

Initially \(G^0=J\), so the input states can all be the same unit
vector. Apply the successive isometries and queries just described.
At the end, factor \(E_0,E_1\) and use one final isometry to produce
their orthogonal direct-sum vectors. Measuring the summand gives
acceptance probability \((E_1)_{x,x}\).
The diagonal equalities in \eqref{ian:pac-query-sdp} preserve unit
norm throughout.

The active channel \(i\) is implemented using response
\(|-\rangle\) at physical index \(i\).
The inert channel is implemented using response \(|+\rangle\)
at any fixed physical index. Thus each phase query costs one
standard coordinate query.

The constructed gates are fixed independently of the unknown
input. Applying them to every \(x\in\Omega_n\) defines the required
global algorithm, including outside \(S\).
\end{proof}

\begin{definition}[General adversary and its dual certificate]
\label{ian:pac-adversary-definition}
For a partial Boolean function \(f:S\to\{0,1\}\), put
\[
 (\Delta_i)_{x,y}=\mathbf1_{\{x_i\ne y_i\}},
\]
and define
\[
 \operatorname{Adv}^{\pm}(f)
  =\max_\Gamma\left\{
       \|\Gamma\|:
       \Gamma=\Gamma^{\mathsf T}\text{ real},\
       \Gamma_{x,y}=0\text{ if }f(x)=f(y),\
       \|\Gamma\circ\Delta_i\|\le1\text{ for all }i
     \right\}.
\]
A dual upper certificate of cost \(A\) consists of positive
semidefinite matrices \(V_i\), indexed by \(S\), with
\[
 \sum_i(V_i)_{x,x}\le A\quad(x\in S),
 \qquad
 \sum_{i:x_i\ne y_i}(V_i)_{x,y}=1
       \quad(f(x)\ne f(y)).
 \tag{32.10}\label{ian:pac-dual-adversary}
\]
Equivalently, factor
\((V_i)_{x,y}=\langle v_{x,i},v_{y,i}\rangle\).
Unlike \eqref{ian:pac-cw}, the constraint is an equality of inner
products, and matrix entries need not be nonnegative.
The matrices may be Hermitian; as in Definition~\ref{ian:pac-query-certificate}, their entrywise real parts are positive semidefinite and satisfy the same constraints, so real certificates suffice.
\end{definition}

\begin{theorem}[Literature adversary statements used here]
\label{ian:pac-adversary-literature}
For a partial Boolean function \(f:S\to\{0,1\}\) on Boolean inputs, write \(Q_\delta(f)\) for the least number of coordinate queries made by a quantum algorithm that outputs \(f(x)\) with probability at least \(1-\delta\) for every \(x\in S\). Then
\[
 Q_\delta(f)\ge
 \frac{1-2\sqrt{\delta(1-\delta)}}2
       \operatorname{Adv}^{\pm}(f)
 \qquad(0\le\delta<1/2).
 \tag{32.11}\label{ian:pac-adversary-lower}
\]
The minimum cost in \eqref{ian:pac-dual-adversary} equals
\(\operatorname{Adv}^{\pm}(f)\).
There is an absolute constant \(C_{\mathrm A}\) such that every
feasible dual certificate of cost \(A\) for a nonconstant \(f\)
constructs, after a constant number of repetitions, an algorithm with error at most \(1/3\) using at most
\(\lceil C_{\mathrm A}A\rceil\) queries.

Consequently, if \(0<\delta<1/2\) and \(r_\delta\) is the least
positive odd integer at least \(18\log(1/\delta)\), a feasible
certificate constructs an acceptance function
\[
 Q\in\mathcal Q_{n,T_\delta},
 \qquad
 T_\delta=r_\delta\lceil C_{\mathrm A}A\rceil,
 \qquad
 |Q(x)-f(x)|\le\delta\quad(x\in S).
 \tag{32.12}\label{ian:pac-adversary-synthesis}
\]
\end{theorem}

\begin{proof}
The lower bound is H{\o}yer--Lee--{\v S}palek, Theorem~2,
Boolean-output case (from literature:
\emph{Negative weights make adversaries stronger},
STOC 2007, 526--535, DOI \(10.1145/1250790.1250867\)).

The dual equality and construction from a feasible solution are
Reichardt, Lemma~2.1 and Section~2 of arXiv:1005.1601v1;
tightness is Theorem~1.3 there (from literature:
\emph{Reflections for quantum query algorithms},
SODA 2011, 560--569, DOI \(10.1137/1.9781611973082.44\)).
That construction, from a feasible solution of cost \(A\), uses \(O(A)\) queries and has constant completeness and soundness, rather than error at most \(1/3\): there are absolute constants \(0\le s<c\le1\) such that it accepts with probability at least \(c\) on one value of \(f\) and at most \(s\) on the other.
Its queries are controlled applications of the phase oracle. A controlled phase query costs one standard coordinate query: route the control-off branch through the inert channel, with response \(|+\rangle\), as in the proof of Theorem~\ref{ian:pac-query-synthesis}.
Repeat the construction independently
\(r_0=\lceil2\ln3/(c-s)^2\rceil\) times and decide by whether the fraction of accepting runs exceeds \((c+s)/2\). Hoeffding's bound gives error at most \(\exp(-r_0(c-s)^2/2)\le1/3\). This defines \(C_{\mathrm A}\).

Independently repeat the resulting algorithm \(r_\delta\) times
and take majority. For each promised input, each answer is correct
with probability at least \(2/3\). The elementary exponential-moment
bound gives majority error at most
\(\exp(-r_\delta/18)\le\delta\).
This proves \eqref{ian:pac-adversary-synthesis}.

For general finite input and output alphabets, the corresponding
bounded-error characterization
\(Q_{1/3}(f)=\Theta(\operatorname{Adv}^{\pm}(f))\)
is Lee--Mittal--Reichardt--{\v S}palek--Szegedy,
Theorem~1.1 (from literature:
\emph{Quantum Query Complexity of State Conversion},
FOCS 2011, 344--353, DOI \(10.1109/FOCS.2011.75\)).
Its queries return symbols; bit encoding must be charged separately.
\end{proof}

\begin{lemma}[Robust compilation into a sector]
\label{ian:pac-robust-compilation}
Let \(Y,Q\in[0,1]\) be defined on a common probability space with
input marginal \(\nu\), and suppose
\[
 \operatorname{Var}Y\ge v,\qquad
 \sup_{F\in\mathcal F_t}|\operatorname{Cov}(Y,F(X))|\le\eta,
 \qquad
 \mathbb E|Q(X)-Y|\le\delta.
\]
Then
\[
 \operatorname{Var}_\nu Q\ge v-2\delta,\qquad
 \sup_{F\in\mathcal F_t}
       |\operatorname{Cov}_\nu(Q,F)|\le\eta+2\delta.
 \tag{32.13}\label{ian:pac-compilation}
\]
If \(Q\in\mathcal Q_{n,T}\) and
\(0<\eta+2\delta<v-2\delta\), this is a realization with
\[
 a=v-2\delta,\qquad \gamma=\eta+2\delta.
\]
Its ambient weight remains \(w\) whenever \(w\nu\le\mu\).
In particular,
\[
 \ell_\mu(t,Q)\ge w(v-\eta-4\delta).
\]
The condition \(4\eta+10\delta\le v\), with \(\gamma>0\),
ensures \(\gamma\le a/4\).
\end{lemma}

\begin{proof}
Let \((Q',Y')\) be an independent copy. Since both differences
\(Q-Q'\) and \(Y-Y'\) belong to \([-1,1]\),
\[
 \left|\operatorname{Var}Q-\operatorname{Var}Y\right|
 =\frac12\left|
      \mathbb E\bigl((Q-Q')^2-(Y-Y')^2\bigr)\right|
 \le\mathbb E|(Q-Y)-(Q'-Y')|\le2\delta.
\]
For \(|F|\le1\),
\[
 |\operatorname{Cov}(Q-Y,F)|
   =|\mathbb E[(Q-Y)(F-\mathbb EF)]|
   \le2\delta.
\]
These prove \eqref{ian:pac-compilation}.
The remaining assertions follow from the sector definition,
Proposition~\ref{ian:hard-sector-forward}, and
\(4(\eta+2\delta)\le v-2\delta\).
\end{proof}

\begin{theorem}[Two adversary-to-sector constructions]
\label{ian:pac-adversary-to-sector}
Let \(f:S\to\{0,1\}\) have a feasible dual certificate of cost
\(A\), and let \(Q\) be the amplified algorithm in
\eqref{ian:pac-adversary-synthesis}.

First, suppose \(\nu(S)=1\), \(w\nu\le\mu\), and
\[
 \operatorname{Var}_\nu f\ge v,\qquad
 \sup_{F\in\mathcal F_t}
       |\operatorname{Cov}_\nu(f,F)|\le\eta.
\]
Whenever \(\eta+4\delta<v\), this constructs a smooth sector
with
\[
 T=T_\delta,\qquad
 a=v-2\delta,\qquad
 \gamma=\eta+2\delta.
\]
If the supplied law is \(\mu(\cdot\mid H)\), it constructs the
corresponding set sector.

Second, suppose \(\mu(S)=1\), \(t\ge1\),
\(0<\epsilon<1/8\), and
\[
 L=\left\lceil512\epsilon^{-5}\right\rceil t,\qquad
 \min_{D\in\mathcal D_L}\mathbb E_\mu|D-f|>\epsilon+\delta.
 \tag{32.14}\label{ian:pac-classical-loss}
\]
Then there is a set \(H\) with membership decided by a tree of
depth below \(L\), such that
\[
 \mu(H)\ge\frac{\epsilon^3}{64},\qquad
 \operatorname{Hard}\left(
 Q,\mu(\cdot\mid H);
 t,\frac{\epsilon^2}{4},\frac{\epsilon^2}{16}
 \right).
 \tag{32.15}
\]
\end{theorem}

\begin{proof}
The first construction follows from
\(\mathbb E_\nu|Q-f|\le\delta\) and
Lemma~\ref{ian:pac-robust-compilation}.

For the second, the triangle inequality gives
\[
 \ell_\mu(L,Q)
 \ge\min_{D\in\mathcal D_L}\mathbb E_\mu|D-f|
       -\mathbb E_\mu|Q-f|>\epsilon.
\]
Apply Corollary~\ref{ian:hard-sector-explicit} in
\path{ian/13_hard_sector_definitions_forward.tex},
with its parameter \(\eta=1/4\).
Its extraction depth is \(L\), weight is
\(\epsilon^3/64\), variance parameter is \(\epsilon^2/4\),
and covariance parameter is \(\epsilon^2/16\).

A distribution satisfying \eqref{ian:pac-classical-loss} can also
be obtained from a sufficiently strong worst-case randomized
classical lower bound by finite minimax. Explicitly, if every
randomized depth-\(L\) tree has some promised input with error
greater than \(\epsilon+\delta\), finite linear programming
duality supplies a law \(\mu\) on \(S\) with that distributional
lower bound for every deterministic Boolean-output tree.
For a Boolean target, allowing leaf values in \([0,1]\) cannot
improve the minimum absolute loss: at each leaf a minimizing
output can be chosen in \(\{0,1\}\).
\end{proof}

\begin{remark}[Scope of the construction and the realization gap]
\label{ian:pac-realization-gap}
A large feasible primal adversary value is a lower bound on the
queries required to decide the specified target. By
\eqref{ian:pac-adversary-lower}, it can reject a proposed query
budget. It supplies neither a small dual certificate nor
law-specific variance, adaptive hiding or ambient weight.
A point-mass sector law has zero variance regardless of the
worst-case adversary value.

A dual certificate guarantees bounded-error decision, rather than
a prescribed acceptance function. It does not imply exact
agreement \(Q=f\), nor agreement with another real target
\(Y\). Amplification supplies the required closeness to \(f\);
Lemma~\ref{ian:pac-robust-compilation} then settles the sector
parameters. For a different \(Y\), one must additionally prove
\(\mathbb E_\nu|Q-Y|\le\delta\), or use
Theorem~\ref{ian:pac-query-synthesis} to certify the intended
acceptance probabilities directly. Exact query-state feasibility
leaves no further finite-dimensional realization gap; discovering
a useful feasible certificate and proving classical hiding remain
the construction problems.
\end{remark}

\begin{proposition}[The weak universal observable converse]
\label{ian:pac-observable-converse}
For every \(Q\in\mathcal Q_{n,T}\), \(0<\epsilon,\rho<1\), repeated
executions estimate \(Q(x)\) to additive error \(\epsilon\), with
failure probability at most \(\rho\), using at most
\[
 T\left\lceil
       \frac{\log(2/\rho)}{2\epsilon^2}
   \right\rceil
\]
coordinate queries.
\end{proposition}

\begin{proof}
Average the independent acceptance bits.
For a Bernoulli variable \(B\),
\(\mathbb E e^{\lambda(B-\mathbb EB)}
 \le e^{\lambda^2/8}\).
For completeness, the logarithm of this moment generating
function has value and derivative zero at zero; its second
derivative is the variance of a tilted Bernoulli law and is at
most \(1/4\). Integrating twice proves the bound.
Exponential Markov bounds for the two tails give error probability
at most \(2e^{-2k\epsilon^2}\) after \(k\) executions.

This gives an estimable observable for every realizable sector,
without identifying an independent mechanism for its algorithm.
Recovering an exact quantization bin additionally requires a
margin from its boundaries.
\end{proof}

\begin{lemma}[Explicit dual certificates from a supplied circuit]
\label{ian:pac-circuit-to-dual}
Suppose a purified \(q\)-query circuit decides
\(f:S\to\{0,1\}\) with error at most
\(0<\delta<1/2\). Its query states explicitly determine a feasible
dual certificate of cost
\[
 A\le\frac{2q}{1-\kappa},
 \qquad \kappa=2\sqrt{\delta(1-\delta)}<1.
 \tag{32.16}\label{ian:pac-circuit-dual-cost}
\]
For an exact circuit the cost bound is \(2q\).
\end{lemma}

\begin{proof}
The argument is the standard dual form of the H{\o}yer--Lee--{\v S}palek lower-bound proof (from literature, cited in the proof of Theorem~\ref{ian:pac-adversary-literature}); we record it because its matrices are explicit.

Let \(\phi_x^r\) be its state immediately before query \(r\),
let \(P_i\) project onto the active phase channel for coordinate
\(i\), and let \(\psi_x\) be its final state.
Define positive semidefinite matrices
\[
 (W_i)_{x,y}
   =2\sum_{r=1}^q
       \langle P_i\phi_x^r,P_i\phi_y^r\rangle,
 \qquad
 G_{x,y}=\langle\psi_x,\psi_y\rangle.
\]
Telescoping the Gram-matrix change of each query gives
\[
 1-G_{x,y}
   =\sum_{i:x_i\ne y_i}(W_i)_{x,y},
 \qquad
 \sum_i(W_i)_{x,x}\le2q.
 \tag{32.17}
\]

Write \(\psi_x=\psi_x^0\oplus\psi_x^1\) for its output components,
and put \(c=(\delta/(1-\delta))^{1/4}\).
Define
\[
 r_x=
 \begin{cases}
 c\psi_x^0\oplus c^{-1}\psi_x^1,&f(x)=0,\\
 c^{-1}\psi_x^0\oplus c\psi_x^1,&f(x)=1.
 \end{cases}
\]
The error bound implies \(\|r_x\|^2\le\kappa\).
If \(f(x)\ne f(y)\), then
\(\langle r_x,r_y\rangle=G_{x,y}\).
Let \(R\) be the Gram matrix of the \(r_x\), and set
\[
 K=\sum_{\ell=0}^\infty R^{\circ\ell},
 \qquad K_{x,y}=\frac1{1-R_{x,y}},
\]
where \(R^{\circ0}=J\).
The series converges entrywise since \(|R_{x,y}|\le\kappa<1\).
Every summand is positive semidefinite, being the Gram matrix of
tensor powers. Hence \(K\) is positive semidefinite and
\(K_{x,x}\le(1-\kappa)^{-1}\).

Set \(V_i=W_i\circ K\). These matrices are positive semidefinite,
again by tensor-product Gram factorization.
For opposite outputs,
\[
 \sum_{i:x_i\ne y_i}(V_i)_{x,y}
   =\frac{1-G_{x,y}}{1-R_{x,y}}=1.
\]
Their diagonal cost satisfies \eqref{ian:pac-circuit-dual-cost}.
Although a convergent series defines \(K\), it is a finite matrix
with the displayed explicit entries and admits a finite Gram
factorization. The matrices \(V_i\) are Hermitian; their entrywise real parts are real positive semidefinite matrices with the same diagonal and the same opposite-output constraint.

For an exact circuit, opposite-output final states are orthogonal.
Thus \(G_{x,y}=0\) on those pairs and \(V_i=W_i\) already satisfies
the dual constraints.
\end{proof}

\begin{proposition}[Character-code certificate and polynomial check]
\label{ian:pac-character-check}
Let \(m\ge4\) be even, \(n=2^m\), and
\[
 c_a(z)=(-1)^{a\cdot z},\qquad
 h(a)=\sum_{j=1}^{m/2}a_{2j-1}a_{2j}\pmod2.
\]
Let \(\nu\) be the uniform law on the character code
\(H=\{c_a:a\in\mathbb F_2^m\}\).
There is an exact one-query algorithm on this promise with global
acceptance function
\[
 Q(x)=\sum_{a:h(a)=1}\alpha_a(x)^2,\qquad
 \alpha_a(x)=\frac1n\sum_zx_z(-1)^{a\cdot z}.
 \tag{32.18}\label{ian:pac-bv-acceptance}
\]
It satisfies
\[
 Q(c_a)=h(a),\qquad
 \operatorname{Var}_\nu Q=\frac{1-1/n}{4},
 \qquad
 \sup_{F\in\mathcal F_t}|\operatorname{Cov}_\nu(Q,F)|
 \le\frac{2^t-1}{2\sqrt n}.
 \tag{32.19}\label{ian:pac-bv-hardness}
\]
In particular, for \(t\le m/2-2\),
\[
 \operatorname{Hard}(Q,\nu;t,3/16,1/8).
\]
A dual certificate and completion weights both have cost \(2\).
The global nonconstant coefficient mass is exactly
\[
 \sum_{S\ne\varnothing}|\widehat Q(S)|
     =\frac{n-1}{2\sqrt n}.
 \tag{32.20}\label{ian:pac-bv-coefficients}
\]
\end{proposition}

\begin{proof}
Prepare \(n^{-1/2}\sum_z|z\rangle\), use phase kickback for one
coordinate query, and apply the Walsh transform.
The resulting amplitudes are \(\alpha_a(x)\).
This is the Bernstein--Vazirani readout (from literature:
E.~Bernstein and U.~Vazirani,
\emph{Quantum Complexity Theory}, SIAM J.~Comput. 26(5),
1411--1473, 1997, DOI \(10.1137/S0097539796300921\)).
The displayed circuit proves the required statement directly.

The Walsh transform is unitary, so
\(\sum_a\alpha_a(x)^2=1\) on the whole cube.
Thus \eqref{ian:pac-bv-acceptance}, together with its complementary
sum, exhibits the degree-one amplitudes and both sums of squares.
On \(c_a\), orthogonality gives the basis state \(|a\rangle\).
This is also an exact query-state certificate in
Theorem~\ref{ian:pac-query-synthesis}.

For explicit dual vectors, use the scalar
\[
 v_{c_a,z}=\sqrt{2/n}.
\]
Distinct characters disagree on exactly \(n/2\) coordinates, so
the opposite-output inner-product sum is one and the cost is two.
Likewise, \(\lambda_z^{c_a}=2/n\) satisfies
\eqref{ian:pac-cw} with cost two.

Put
\[
 b(s)=\frac1n\sum_a h(a)(-1)^{a\cdot s}.
\]
The Walsh sum of \((-1)^{h(a)}\) factors over its two-bit blocks.
Each normalized block sum has absolute value \(1/2\).
Consequently,
\[
 b(0)=\frac{1-n^{-1/2}}2,\qquad
 |b(s)|=\frac1{2\sqrt n}\quad(s\ne0).
\]
This proves the variance formula.
A classical coordinate query to \(c_a(z)\) reveals one linear
equation in \(a\). Each positive-probability adaptive leaf
therefore conditions \(a\) onto an affine subspace of codimension
at most \(t\). Its conditional mean of \(h-b(0)\) is a sum of at
most \(2^t-1\) nonconstant Fourier terms, each of absolute value
\(1/(2\sqrt n)\). Averaging against arbitrary leaf outputs in
\([-1,1]\) proves \eqref{ian:pac-bv-hardness}.
The stated sector parameters follow.

Expanding the global acceptance polynomial gives constant
coefficient \(b(0)\), and, for unordered distinct coordinates
\(z,w\),
\[
 \widehat Q(\{z,w\})=\frac2n\,b(z+w).
\]
All its other nonconstant coefficients vanish. There are
\(n(n-1)/2\) such pairs, proving
\eqref{ian:pac-bv-coefficients}.
Moreover,
\[
 \operatorname{Cov}_\nu(Q,x_zx_w)=b(z+w),
 \qquad
 \operatorname{Var}_\nu Q
   =\frac1{2\sqrt n}
       \sum_{S\ne\varnothing}|\widehat Q(S)|.
\]
Thus \eqref{ian:pac-pm-display} is verified whenever \(t\ge2\):
hardness then forces \(\gamma\ge1/(2\sqrt n)\).

Under ambient law \(\mu=\nu\), the sector has weight one.
Under the uniform cube law, its weight is
\(n2^{-n}\), and
\(\nu\le(2^n/n)U_n\).
These are the corresponding parameters for
Theorem~\ref{ian:pac-product-exclusion}.
\end{proof}

\begin{proposition}[Simon in physical coordinate form]
\label{ian:pac-simon-check}
Let \(m\ge10\), \(M=2^m\), \(G=\mathbb F_2^m\), and \(n=mM\).
Choose \(s\) uniformly from \(G\setminus\{0\}\), assign distinct
uniformly random labels in \(G\) to the \(M/2\) pairs
\(\{u,u+s\}\), and denote the resulting table by \(f_X:G\to G\).
Encode it as
\[
 X_{u,j}=(-1)^{f_X(u)_j}.
\]
Let \(\nu\) be this promise law, \(H=\operatorname{supp}\nu\),
and \(Y=s_1\).

For \(k\ge m-1\), there is a globally defined algorithm using
\(T=mk\) coordinate queries with
\[
 Q(X)=p_{m,k}Y(X)\quad(X\in H),
 \qquad
 p_{m,k}=\prod_{j=0}^{m-2}(1-2^{j-k}),
 \qquad
 1-p_{m,k}\le2^{m-1-k}.
 \tag{32.21}\label{ian:pac-simon-acceptance}
\]
For \(c_t=\binom t2<M-1\), its covariance satisfies
\[
 \sup_{F\in\mathcal F_t}|\operatorname{Cov}_\nu(Q,F)|
 \le p_{m,k}\frac{2c_t}{M-1-c_t}.
 \tag{32.22}\label{ian:pac-simon-hiding}
\]
Taking
\[
 k=m+9,\qquad
 t=\left\lfloor\frac{\sqrt M}{32}\right\rfloor
\]
gives a sector of weight one under \(\mu=\nu\), with
\[
 T=m(m+9),\qquad a=\frac18,\qquad\gamma=\frac1{64}.
\]
Both an explicit query-state certificate and an explicit dual
certificate of cost less than \(3T\) are available.
\end{proposition}

\begin{proof}
For every cube input define
\(f_X(u)_j=(1-X_{u,j})/2\).
Its whole-output XOR oracle uses exactly \(m\) coordinate queries,
one for each output bit. Start each sample with
\[
 M^{-1/2}\sum_{u\in G}|u\rangle|0^m\rangle,
\]
apply that oracle, apply the Walsh transform to the first
register, and retain both registers. No inverse oracle call is
required.

This is Simon's Fourier-sampling procedure (from literature:
D.~R.~Simon, \emph{On the Power of Quantum Computation},
SIAM J.~Comput. 26(5), 1474--1483, 1997,
DOI \(10.1137/S0097539796298637\)); the calculations needed here
follow directly.

Indeed, the global amplitude of \(|v,z\rangle\) in one sample is
\[
 A_{v,z}(X)=\frac1M\sum_{u\in G}(-1)^{u\cdot v}
       \prod_{j=1}^m\frac{1+(-1)^{z_j}X_{u,j}}2.
 \tag{32.23}
\]
It has degree at most \(m\).
The sample probabilities are
\[
 B_v(X)=\sum_z A_{v,z}(X)^2
 =\frac1{M^2}\sum_{u,w\in G}(-1)^{(u+w)\cdot v}
       \prod_{j=1}^m\frac{1+X_{u,j}X_{w,j}}2.
 \tag{32.24}
\]
They are nonnegative and sum to one for every cube input.

Take \(k\) independent samples. If their vectors
\(v_1,\ldots,v_k\) have rank \(m-1\), output the first bit of the
unique nonzero vector in their common orthogonal complement.
At every other rank, output zero.
This is a fixed rule on all cube inputs.

Let \(\mathcal A\) be the set of sample tuples for which that rule
outputs one. The global acceptance function is
\[
 Q(X)=\sum_{(v_1,\ldots,v_k)\in\mathcal A}
          \prod_{r=1}^k B_{v_r}(X).
 \tag{32.25}
\]
Equivalently, it is the sum of squares of the amplitude products
\(\prod_r A_{v_r,z_r}(X)\) over accepting basis states.
The complementary basis states give \(1-Q\).
These formulas, together with the stated sequence of coordinate
queries, explicitly certify query compatibility, amplitude degree
at most \(mk\), acceptance degree at most \(2mk\), and global
boundedness.

On the promise, equal outputs occur precisely when
\(w=u\) or \(w=u+s\). Hence
\[
 B_v(X)=\frac{1+(-1)^{s\cdot v}}M
       =\frac2M\mathbf1_{\{v\in s^\perp\}}.
\]
The samples are independent uniform vectors in the
\((m-1)\)-dimensional space \(s^\perp\).
The full-rank probability is
\(\prod_{j=0}^{m-2}(1-2^{j-k})\).
Alternatively, failure implies that one of the
\(2^{m-1}-1\) nonzero linear functionals vanishes on all samples,
so its probability is at most \(2^{m-1-k}\).
This proves \eqref{ian:pac-simon-acceptance}, including the
specified failure output.

We next prove adaptive hiding.
A deterministic whole-output strategy may cache repeated arguments.
For a collision-free transcript with distinct queried arguments
\(u_1,\ldots,u_q\), put
\[
 E_\tau=\{u_i+u_j:i<j\}.
\]
For any candidate nonzero mask \(s'\notin E_\tau\), these arguments
lie in distinct \(s'\)-pairs. The likelihood of their prescribed
distinct output labels is exactly \(1/(M)_q\), independently of
\(s'\), where \((M)_j=M(M-1)\cdots(M-j+1)\) denotes the falling factorial. For \(s'\in E_\tau\) it is zero.
Adaptive query choices are already determined by the preceding
answers in the fixed transcript. Therefore the posterior mask is
uniform on
\[
 (G\setminus\{0\})\setminus E_\tau.
\]

Write
\[
 \rho=\mathbb E_\nu Y=\frac{M}{2(M-1)}.
\]
At a collision-free terminal transcript,
\[
 |\mathbb E[Y\mid\tau]-\rho|
 \le\frac{|E_\tau|}{M-1-|E_\tau|}
 \le\frac{c_t}{M-1-c_t}.
\]
After \(j\) collision-free queries, a new argument can collide with
at most \(j\) previous arguments, so its conditional collision
probability is at most \(j/(M-1-c_t)\).
Summing over first-collision steps bounds the probability of any
collision by \(c_t/(M-1-c_t)\).
On collision transcripts the posterior deviation is at most one.
Consequently,
\[
 \mathbb E|\mathbb E[Y\mid\tau]-\rho|
 \le\frac{2c_t}{M-1-c_t}.
\]
A depth-\(t\) bit tree can be simulated by at most \(t\)
whole-output queries, caching each returned symbol.
Its bit transcript is a function of the simulated transcript;
conditional Jensen preserves the bound.
Multiplying the posterior deviation by a leaf output of absolute
value at most one proves the covariance bound for \(Y\).
Since \(Q=p_{m,k}Y\) on the promise, this proves
\eqref{ian:pac-simon-hiding}.

The target variance is
\[
 \operatorname{Var}_\nu Y
   =\frac14-\frac1{4(M-1)^2}.
\]
For \(k=m+9\), \(p_{m,k}\ge1-2^{-10}\), so
\[
 \operatorname{Var}_\nu Q
   =p_{m,k}^2
       \left(\frac14-\frac1{4(M-1)^2}\right)>\frac18.
\]
For the stated \(t\), \(c_t\le M/2048\).
As \(M\ge1024\),
\[
 \frac{2c_t}{M-1-c_t}
 \le\frac2{2045}<\frac1{512}<\frac1{64}.
\]
The announced hard-sector parameters follow.

To exhibit dual matrices, take the actual prequery and final
states of this circuit, specified by the sample preparation and
rank postprocessing above, and apply the explicit formulas
\(W_i,R,K,V_i\) in Lemma~\ref{ian:pac-circuit-to-dual}, with
\(\delta=2^{-10}\).
They satisfy \eqref{ian:pac-dual-adversary} and have cost
\[
 A\le
 \frac{2T}{1-2\sqrt{2^{-10}(1-2^{-10})}}<3T.
\]
This supplies finite matrices with specified entries, rather than
an inference from completion weights.

Finally, every table uniquely determines its nonzero mask, and
each mask has \((M)_{M/2}\) injective pair labelings. Thus
\[
 |H|=(M-1)(M)_{M/2},\qquad
 U_n(H)=\frac{(M-1)(M)_{M/2}}{M^M}.
\]
The promise law is uniform on \(H\), so
\[
 \nu\le
 \frac{M^M}{(M-1)(M)_{M/2}}\,U_n.
\]
The weight is one under the promise law and is the displayed
small probability under the uniform product law.
\end{proof}

\begin{remark}[Simon and the coefficient obstruction]
\label{ian:pac-simon-pm}
For Simon, physical degree is at most \(2T=2mk\), even though
\(Q|_H=p_{m,k}s_1\) is linear in one hidden mask bit.
If \(t<2T\), the applicable necessary statement is
\eqref{ian:pac-pm-tail}; \eqref{ian:pac-pm-display} must not be invoked without its
depth hypothesis.

If \(t\ge2T\), as occurs for all sufficiently large \(m\) in the
stated package, \eqref{ian:pac-pm-display} gives
\[
 \sum_{S\ne\varnothing}|\widehat Q(S)|\ge8.
\]
A sharper check follows by applying the proved Simon hiding bound
at depth \(2T\). Write \(c=\binom{2T}{2}\), and suppose \(c<M-1\).
Every nonconstant character in the global acceptance polynomial
then has covariance with \(Q\) at most
\(p_{m,k}\,2c/(M-1-c)\). Consequently,
\[
 \sum_{S\ne\varnothing}|\widehat Q(S)|
 \ge
 \frac{
 p_{m,k}\operatorname{Var}_\nu Y\,(M-1-c)
 }{2c}.
\]
Thus the explicit circuit has the coefficient mass forced by
adaptive hiding. Its amplitude construction, not a polynomial
specified only on \(H\), certifies its realizability.
\end{remark}
